 \documentclass{article}
\usepackage{amsmath}
\usepackage{graphicx}
\graphicspath{ {images/} }
\usepackage[spanish]{babel}
\usepackage{bbm}
\usepackage{amsthm}
\usepackage{authblk}
\usepackage{hyperref}
\usepackage[utf8]{inputenc}
\usepackage{mathrsfs}
\usepackage{amsfonts}
\usepackage{amssymb}
\usepackage{enumerate}
\usepackage[left=3cm,right=2cm,top=2cm,bottom=2cm]{geometry}
\usepackage{epsfig}
 \renewcommand{\baselinestretch}{0.93}
 \thispagestyle{empty}
 \pagestyle{empty}
\setlength{\textheight}{53pc} \setlength{\textwidth} {36pc}
\setlength{\topmargin}{-4mm}
\usepackage{multicol}
\newcommand*{\email}[1]{%
    \normalsize\href{mailto:#1}{#1}\par
    }
\setlength{\columnsep}{1cm}
\newcommand{\N}{\mathbb{N}}
\newcommand{\calM}{\cal{M}}
\newcommand{\calP}{\cal{P}}
\newcommand{\F}{\mathbb{F}}
\newcommand{\R}{\mathbb{R}}
\newcommand{\Z}{\mathbb{Z}}
\newcommand{\Q}{\mathbb{Q}}
\newcommand{\C}{\mathbb{C}}
\newcommand{\bbH}{\mathbb{H}}


\theoremstyle{definition}
\newtheorem{ejer}{Ejercicio}
\author{Marcos Morales}
\affil{Universidad Finis Terrae}
\title{Primera Semana\\}



	


\begin{document}


\begin{picture}(400, 75)
   \put(-15,15){\includegraphics[width=5cm]{UFT.png}}
\end{picture}



\begin{center}
\textbf{\large{Laboratorio Evaluado 2}} \\
\end{center}

\setlength{\parindent}{0cm}



\begin{center}
\textbf{Cálculo Vectorial}
\end{center}


\begin{center}
\textbf{\large{Análisis del Modelo}} \\
\end{center}



\hspace{100cm} \\ 


Considere que los dos modelos Cobb-Douglas de dos empresas $A$ y  $B$ están dados por

\begin{align*}
    Y_A &= 10K^{0.3}L^{0.7} \\
    Y_B &= 7K^{0.83}L^{0.17}
\end{align*}

El objetivo de este laboratorio es que se realize un análisis exaustivo del modelo Cobb-Douglas, comparando las producciones de estas dos empresas. \\

\begin{center}
    \textbf{\Large{Actividades}}
\end{center}

\begin{itemize}
    \item Usando Matlab u Ocatve, realice un gráfico de ambos modelos.
    \item Los modelos presentan máximos relativos?, mínimos relativos?, puntos silla?. Argumente usando el gradiente de la función. 
    \item Compare los modelos en distintos puntos del plano, incluya al menos 10 puntos.
    \item Considere que el factor capital y la dotación del trabajo están relacionados por \\$K+L=C$, donde $C$ es una constante. Determine  los puntos extremos con esa condición
    \item ¿Existe algún valor de $C$ para el cual el máximo relativo del punto anterior de $Y_A$ sea mayor al de $Y_B$ ?. En caso de ser cierto determine los valores posibles.
    \item Explique la razón matemática por la cuál ocurre esto 
    \item Considere ahora $Y_D = 7K^{0.5}L^{0.5}$ e $Y_E = 7K^{0.5}L^{0.5}$, y realice la comparación de los máximos relativos de $Y_A$, $Y_B$, $Y_D$ e $Y_E$
\end{itemize}









\newpage


\textbf{Solución: } 1) No pude instalar matlab en mi Pc pero los gráficos son los estandar de un modelo Cobb-Douglas. \\



2) Calculamos el gradiente de $Y_A$, tenemos

\begin{align*}
    \nabla Y_A = (3K^{-0,7}L^{0,7},7K^{0,3}L^{-0,3})
\end{align*}

Necesitamos entonces que

\begin{align*}
    3K^{-0,7}L^{0,7} &=0 \\
    7K^{0,3}L^{-0,3} &=0
\end{align*}

Este sistema no tiene solución pues $K \neq 0$ y $L \neq 0$, porque no se puede dividir por cero. Por lo tanto no hay puntos críticos, para $Y_B$ es análogo.\\ \\



2)Los puntos que se pueden elegir son muchos solo hay que comprobar que sean consistentes, por ejemplo en el punto $(1,1)$, tenemos $Y_A=10$ e $Y_B=7$ por lo tanto $Y_A(1,1) > Y_B(1,1)$. \\ \\


3)Tenemos en cada caso la restricción $g(K,L)=K+L-C$, usamos la regla de multiplicadores de Lagrange, necesitamos 

\begin{align*}
    \nabla Y_A &= \lambda \nabla g \\
     (3K^{-0,7}L^{0,7},7K^{0,3}L^{-0,3}) &= \lambda (1,1)
\end{align*}

Luego tenemos el sistema
 
\begin{align*}
    3K^{-0,7}L^{0,7} &= \lambda \\
    7K^{0,3}L^{-0,3} &= \lambda
\end{align*}

Entonces tenemso que


\begin{align*}
    3K^{-0,7}L^{0,7} &= 7K^{0,3}L^{-0,3}\\
    3L &= 7K \\
    L&= \frac{7K}{3}
\end{align*}

Reemplazando en $K+C=L$, obtenemos

\begin{align*}
    K+ \frac{7K}{3} &= C \\
    K &= \frac{3C}{10}
\end{align*}

y entonces 

\begin{align*}
    L&= \frac{7C}{10}
\end{align*}

El punto extremo para $Y_A$ es $\displaystyle \left( \frac{3C}{10}, \frac{7C}{10}\right)$. \\ 

Si hacemos los mismo pasaos para $Y_B$, obtenemos el punto obtenemos el punto extremo condicionado   $\displaystyle \left( \frac{83C}{100}, \frac{17C}{100}\right)$ \\


Si evaluamos los puntos encontrados en $Y_A$ e $Y_B$, obtenemos

\begin{align*}
    Y_A &= 10(0,3C)^{0,3}(0,7C)^{0,7} = 10(0,3)^{0,3}(0,7)^{0,7}C \sim 5,4C \\
    Y_B &= 10(0,83C)^{0,83}(0,17C)^{0,17} = 10(0,83)^{0,83}(0,17)^{0,17}C \sim 4,4C\\
\end{align*}

4) Podemos ver que si $C$ es mayor que $0$ siempre sew tendrá que el máximo de $Y_A$ será más grnade independiente del valor de $C$. \\\\


5) La razón por la cual esto ocurre es que en general el máximo de una Cobbdouglas de la forma $Y=AK^{\alpha}L^{1-\alpha}$ bajo la restricción $K+L=C$ es

\begin{align*}
    (\alpha C, (1-\alpha)C )
\end{align*}

Como se puede inferir y demostrar, al reemplzar en la función obtenemos

\begin{align*}
    Y = AC\alpha^\alpha(1-\alpha)^{1-\alpha}
\end{align*}

Por lo tanto depende de la función $\alpha^\alpha(1-\alpha)^{1-\alpha}$ y del valor de $A$, si considereamos

\begin{align*}
    f(\alpha) =  \alpha^\alpha(1-\alpha)^{1-\alpha}
\end{align*}

Entonces

\begin{align*}
    f'(\alpha) =  \alpha^\alpha(1-\alpha)^{1-\alpha} \ln\left( \frac{\alpha}{1-\alpha}\right)
\end{align*}

Podemos ver entonces qye esta función tiene derivada positiva si $\alpha > 0,5$ y negativa si $\alpha < 0,5$, por lo tanto es decreciente entre $[0,1/2]$ y creciente en $[1/2,1]$ y además es simétrica respecto a $1/2$. Aí que en un principio $Y_B$ debería ser mayor que $Y_A$, pero por el facotor de ajuste $10$ que es mayor que $7$ al final termina ocurriendo que $Y_B$ es más pequeño. \\ \\

6) Primero que nada tenemos que $Y_D=Y_E$ (esto es solo una distracción), como pudimos ver en el punto 5 el valor más bajo se va a alcanzar cuando $\alpha=0,5$, por lo tanto como tiene el mismo factor $A=7$ que $Y_B$, podemos concluir directamente que en la evaluación de los puntos máximos se tendrá que

\begin{align*}
    Y_D=Y_E < Y_B < Y_A
\end{align*}



\end{document}